\documentclass[10pt,twocolumn]{article}

% ── Packages ──────────────────────────────────────────────────────────
\usepackage[utf8]{inputenc}
\usepackage[T1]{fontenc}
\usepackage{lmodern}
\usepackage[margin=0.75in]{geometry}
\usepackage{booktabs}
\usepackage{enumitem}
\usepackage{hyperref}
\usepackage{xcolor}
\usepackage{graphicx}
\usepackage{amsmath,amssymb}
\usepackage{natbib}
\usepackage{caption}
\usepackage{listings}
\usepackage{float}

\lstset{
  basicstyle=\scriptsize\ttfamily,
  breaklines=true,
  columns=fullflexible,
  keepspaces=true,
  showstringspaces=false,
  xleftmargin=0.5em,
  xrightmargin=0.5em,
  frame=single,
  rulecolor=\color{gray!40},
  aboveskip=4pt,
  belowskip=4pt,
}
\usepackage{microtype}
\usepackage{titlesec}

\hypersetup{
  colorlinks=true,
  linkcolor=blue!60!black,
  citecolor=blue!60!black,
  urlcolor=blue!60!black,
}

\titlespacing*{\section}{0pt}{1.2ex plus .3ex minus .2ex}{0.6ex plus .1ex}
\titlespacing*{\subsection}{0pt}{0.9ex plus .2ex minus .1ex}{0.4ex plus .1ex}

% ── Title ─────────────────────────────────────────────────────────────
\title{\textbf{EABench: A Configurable Benchmark Framework\\for Evaluating LLM-Powered Enterprise Agents}}

\author{
  EABench Contributors\\
  \texttt{https://github.com/xhyccc/EABench}
}

\date{}

\begin{document}
\maketitle

% ══════════════════════════════════════════════════════════════════════
% ABSTRACT
% ══════════════════════════════════════════════════════════════════════
\begin{abstract}
Large language model (LLM) agents are increasingly deployed in enterprise
productivity settings---automatically managing emails, scheduling meetings,
retrieving files, and executing multi-step workflows across heterogeneous data
sources.  Despite rapid progress in general-purpose agent benchmarks, existing
evaluation frameworks fall short in two important respects: they provide little
support for generating \emph{realistic enterprise data at scale}, and they
offer limited flexibility in \emph{agent architecture configuration}.  We
present \textbf{EABench}, an open-source benchmark framework designed to
address both gaps.  EABench provides (1)~an LLM-driven \emph{synthetic data
generation} pipeline that produces coherent, inter-connected corporate corpora
spanning emails, meetings, files, chats, and channels; (2)~a
\emph{configurable agent runtime} that supports multiple execution strategies
(ReAct and Plan-and-Execute) and multiple LLM backends through a
provider-agnostic interface; (3)~a comprehensive \emph{enterprise tool suite}
encompassing semantic search over all content types, sandboxed file I/O, and
code execution; and (4)~a \emph{multi-dimensional evaluation harness} combining
deterministic assertions, citation relevance scoring, and LLM-as-a-Judge for
end-to-end and user-level evaluation.  Together, these components enable
systematic, reproducible comparison of agent designs under realistic enterprise
conditions.
\end{abstract}

% ══════════════════════════════════════════════════════════════════════
% 1  INTRODUCTION
% ══════════════════════════════════════════════════════════════════════
\section{Introduction}
\label{sec:intro}

The emergence of LLM-based autonomous agents capable of using tools, planning
multi-step workflows, and interacting with external APIs has prompted a wave of
benchmarks aimed at measuring agent capability.  However, a growing gap exists
between academic benchmark environments and the conditions found in real
enterprise deployments.  Enterprise AI agents must navigate large,
\emph{heterogeneous} corpora---email threads, calendar invites, document
repositories, instant messages, and project channels---where information is
\emph{fragmented} across sources, \emph{access-controlled} at the user level,
and \emph{temporally distributed} across weeks or months of organizational
history.

Current benchmarks such as WebArena~\citep{zhou2023webarena},
AgentBench~\citep{liu2023agentbench}, and ToolBench~\citep{qin2023toolllm}
focus primarily on web navigation, software engineering tasks, and API calling,
respectively.  While these settings surface important capabilities, they do not
capture the \emph{enterprise information retrieval and synthesis} tasks that
constitute the bulk of knowledge-worker workload: ``Summarize all action items
from last week's incident post-mortems,'' ``Find the latest contract
negotiation emails with Vendor~X,'' or ``Who owns the on-call rotation this
month?''

EABench directly targets this gap.  Its three core contributions are:

\begin{enumerate}[leftmargin=*,itemsep=2pt]
  \item \textbf{Realistic, configurable data generation.}
    EABench uses an LLM-driven pipeline to generate fully synthetic yet
    coherent enterprise tenants at any scale, with user-defined industry,
    company size, timeline, and storyline events.  The resulting corpora are
    inter-connected: emails reference meeting decisions, files document project
    outcomes, and chats discuss ongoing incidents---exactly as in real
    organizations.

  \item \textbf{Flexible agent architecture.}
    Rather than fixing a single agent design, EABench exposes a
    \emph{configuration-as-code} model in which every aspect of an
    agent---LLM backend, execution strategy (ReAct vs.\ Plan-and-Execute),
    available tools, system prompt, and planning prompt---is specified in a
    human-readable YAML file.  This enables controlled ablation studies and
    rapid iteration without modifying source code.

  \item \textbf{Comprehensive evaluation.}
    EABench's evaluation harness combines \emph{deterministic assertions}
    (verifiable ground-truth checks), \emph{citation scoring} (verifying that
    responses cite real, relevant source entities), and
    \emph{LLM-as-a-Judge} scoring for qualitative dimensions such as
    faithfulness, completeness, and reasoning quality.  User-level evaluation
    simulates realistic access-control scenarios by running queries under
    specific user identities that determine data visibility.
\end{enumerate}

The remainder of this paper is organized as follows.
Section~\ref{sec:related} surveys related work.
Section~\ref{sec:method} describes EABench's methodology in detail.
Section~\ref{sec:experiments} presents experimental results and analysis.
Section~\ref{sec:conclusion} concludes with a discussion of limitations and
future directions.

% ══════════════════════════════════════════════════════════════════════
% 2  RELATED WORK
% ══════════════════════════════════════════════════════════════════════
\section{Related Work}
\label{sec:related}

\subsection{Agent Benchmarks}

\textbf{WebArena}~\citep{zhou2023webarena} provides a realistic web browsing
environment with four web applications (shopping, reddit, GitLab, map) and 812
long-horizon tasks.  While WebArena evaluates task completion in a realistic
setting, its tasks are scoped to web interfaces and do not address enterprise
knowledge retrieval.  \textbf{WorkArena}~\citep{drouin2024workarena} extends
this approach to ServiceNow workflows, targeting enterprise service management
tasks.  However, its data is limited to a single application and does not cover
the richly inter-connected multi-source data of a real enterprise.

\textbf{AgentBench}~\citep{liu2023agentbench} is a multi-dimensional benchmark
spanning eight environments including OS, database, web browsing, and lateral
reasoning tasks.  It demonstrates that current LLMs struggle with multi-step
tool-use tasks; however, its environments are largely synthetic or game-like
rather than representative of knowledge-worker contexts.

\textbf{ToolBench}~\citep{qin2023toolllm} focuses specifically on evaluating
LLMs' ability to invoke real-world APIs from a large catalog of 16,000 APIs.
While this directly tests tool-use, the tasks are API-centric rather than
information-synthesis-centric, and there is no notion of user-level access
control or data provenance.

\textbf{$\tau$-bench}~\citep{yao2024taubench} introduces the concept of user
simulation for interactive task evaluation in airline and retail domains.  The
approach of simulating realistic user interactions---rather than specifying
fully deterministic tasks---is conceptually related to EABench's user-level
evaluation mode, though $\tau$-bench targets customer service scenarios rather
than enterprise knowledge work.

\subsection{Enterprise AI Assistants}

Several industry systems have been deployed for enterprise AI assistance.
\textbf{Microsoft 365 Copilot}~\citep{microsoft2023copilot} integrates LLMs
with Microsoft Graph to answer queries over emails, documents, meetings, and
Teams chats.  Its architecture involves retrieval-augmented generation (RAG)
over enterprise data with user-level access control.  While such systems
demonstrate the commercial viability of enterprise agents, they do not provide
open evaluation datasets or benchmarking infrastructure.

\textbf{Glean} and similar enterprise search platforms combine vector search
with access control lists (ACLs) to implement personalized search across
connected SaaS applications.  These systems illustrate the importance of
user-context-aware retrieval---a feature EABench explicitly supports through
its user-level indexing and filtering.

\subsection{Synthetic Data Generation for Benchmarks}

Several recent works have investigated LLM-generated synthetic data for
evaluation.
\textbf{AgentInstruct}~\citep{mitra2024agentinstruct} uses LLMs to generate
instruction tuning data by transforming raw documents through a pipeline of
multiple specialized agents.
\textbf{FrontierMath}~\citep{glazer2024frontiermath} demonstrates that
LLM-generated problems can challenge state-of-the-art models.
\textbf{WildChat}~\citep{zhao2024wildchat} collects naturally occurring LLM
conversations in the wild.

EABench's approach differs from prior work in that it generates
\emph{structured organizational corpora} rather than isolated question-answer
pairs.  The generator produces inter-connected narrative artifacts (emails,
meetings, files, chats) grounded in a shared storyline, enabling queries that
require multi-hop reasoning across document types and temporal reasoning over
event sequences.

\subsection{LLM-as-a-Judge Evaluation}

Automated evaluation using LLMs as judges has gained significant traction
following MT-Bench~\citep{zheng2023judging} and Chatbot
Arena~\citep{chiang2024chatbot}.  LLM judges show high correlation with human
preferences at lower cost and enable evaluation of open-ended responses that
resist automatic metrics.

EABench extends the LLM-as-a-Judge paradigm to \emph{agentic evaluation},
where the judge must assess not only the final response but also the
\emph{reasoning process}---including whether cited sources are accurate,
whether tool calls were appropriate, and whether the overall research plan was
followed.

\subsection{ReAct and Plan-and-Execute Agent Frameworks}

\textbf{ReAct}~\citep{yao2022react} introduced the Reason+Act paradigm, where
LLMs interleave reasoning traces with tool invocations.  This approach
demonstrated that explicitly reasoning about the next action before taking it
substantially improves task performance on knowledge-intensive tasks.

\textbf{Plan-and-Solve}~\citep{wang2023plansolve} agents extend this idea by
separating the planning phase---generating a high-level task decomposition---from
the execution phase.  This separation can improve reliability on complex
multi-step tasks by preventing the agent from ``losing the thread'' during long
tool-use sequences.

EABench implements both paradigms as first-class, YAML-configurable execution
strategies, enabling direct comparison of their performance characteristics
across benchmark tasks.

% ══════════════════════════════════════════════════════════════════════
% 3  METHODOLOGY
% ══════════════════════════════════════════════════════════════════════
\section{Methodology}
\label{sec:method}

EABench is organized around four interacting subsystems:
(1)~the data generation pipeline,
(2)~the configurable agent runtime,
(3)~the enterprise tool suite, and
(4)~the evaluation harness.
Figure~\ref{fig:arch} shows the overall system architecture.

\begin{figure*}[t]
\centering
\small
\begin{verbatim}
+----------------------------------------------------------------------+
|                     DATA GENERATION PIPELINE                         |
|  Story Config --> LLM-driven generation --> Tenant YAML + Files      |
+-----------------------------+----------------------------------------+
                              | tenant data
                              v
+----------------------------------------------------------------------+
|                      AGENT RUNTIME LAYER                             |
|                                                                      |
|  +-----------------+  +------------------+  +-------------------+    |
|  | Agent Config    |  |  LLM Provider    |  |   Sandbox         |    |
|  | (YAML)          |  | OpenAI / Azure   |  |  LocalSandbox     |    |
|  |  - model        |  | Anthropic / Local|  |  DockerSandbox    |    |
|  |  - flow strategy|  +--------+---------+  +-------+-----------+    |
|  |  - tools        |           |                    |                |
|  |  - prompts      |  +--------v--------------------v-----------+    |
|  +-----------------+  |          AgentRunner                    |    |
|                       |  ReAct loop / Researcher strategy       |    |
|                       +--------------------+--------------------+    |
|                                            |                         |
|  +-----------------------------------------v----------------------+  |
|  |                      Tool Suite                                |  |
|  | search_email  search_file  search_meeting  search_people  ...  |  |
|  | read_file  list_files  execute_command  execute_python         |  |
|  +----------------------------+-----------------------------------+  |
|                               |                                      |
|  +----------------------------v-----------------------------------+  |
|  |                   Search Engine                                |  |
|  |    Vector indices (per content type) + ACL filtering           |  |
|  +----------------------------------------------------------------+  |
+-----------------------------+----------------------------------------+
                              | agent response + traces
                              v
+----------------------------------------------------------------------+
|                      EVALUATION HARNESS                              |
| Assertions | Citation Scoring | LLM-as-a-Judge | A/B Compare        |
+----------------------------------------------------------------------+
\end{verbatim}
\caption{EABench system architecture.  The data generation pipeline produces
tenant corpora consumed by the agent runtime.  The agent uses LLM providers,
tools, and a search engine to answer queries inside a sandbox.  The evaluation
harness scores agent responses along multiple dimensions.}
\label{fig:arch}
\end{figure*}

\subsection{Data Generation}
\label{sec:datagen}

\textbf{Design goals.}
Benchmark data for enterprise agents must satisfy three requirements
simultaneously:
(a)~\emph{realism}---the content must resemble actual corporate communication
artifacts;
(b)~\emph{coherence}---distinct artifacts (emails, meetings, files) must form a
consistent narrative rather than isolated documents; and
(c)~\emph{controllability}---the researcher must be able to specify the domain,
scale, and scenario characteristics.

\textbf{Generator architecture.}
The data generation pipeline is driven by a \texttt{DataGenerator} class that
accepts a \texttt{StoryConfig} specifying the company name, industry, size,
number of users, timeline length, and a list of named \emph{story events}
(e.g., ``Product Alpha Launch,'' ``Production Database Outage'').  The
generator proceeds in five phases:

\begin{enumerate}[leftmargin=*,itemsep=2pt]
  \item \textbf{User generation.}
    Synthetic employee profiles are created with realistic names, email
    addresses, departments, job titles, management hierarchies, and skill sets.
    Users are assigned to organizational groups that determine their
    communication patterns.

  \item \textbf{Story arc construction.}
    For each story event, the generator synthesizes a brief scenario narrative
    that will inform all subsequent content generation.  This ensures that
    emails, meetings, and files reference the same shared context.

  \item \textbf{Daily content generation.}
    For each day in the specified timeline, the generator prompts the LLM to
    produce a batch of emails, meeting records, instant messages, and file
    documents consistent with the day's story state.  To maintain coherence
    across days, the generator maintains a rolling summary of past events.

  \item \textbf{Cross-referencing.}
    The generator explicitly instructs the LLM to create cross-references:
    emails should cite decisions from meetings, files should document outcomes
    of email discussions, and chats should reference both.

  \item \textbf{Evaluation query generation.}
    Given a completed tenant, a separate pipeline generates evaluation queries
    grounded in the tenant's content, with LLM-written assertion descriptions
    that a Judge can later verify.
\end{enumerate}

\textbf{Example data-generator configuration.}
Listing~\ref{lst:storyconfig} shows a representative \texttt{StoryConfig}.
The \texttt{key\_events} field provides the narrative arc that drives
cross-document coherence: every email, meeting, and file generated for
the tenant will reference one or more of these events.
The generator produces fully interconnected corpora---emails discussing
launch readiness point to the same meeting transcripts in which
the go/no-go decision was made, and incident post-mortems in the file
store cite the chat threads in which the outage was first flagged.

\begin{lstlisting}[caption={Example StoryConfig for a software-startup tenant.},label={lst:storyconfig}]
# Data-generator configuration (StoryConfig)
company_name: TechNova Inc.
industry: Software-as-a-Service
company_size: small
num_users: 5
duration_days: 3
key_events:
  - "Product Alpha Launch: TechNova releases its flagship SaaS
     platform to enterprise customers after months of testing"
  - "Production Outage: An untested code push crashes the live
     API; on-call engineers execute an emergency rollback"
description: >
  A fast-growing SaaS startup building enterprise productivity
  tools for mid-market companies, now undergoing a critical
  product-launch phase.
\end{lstlisting}

\textbf{Strengths over prior approaches.}
Unlike BEIR~\citep{thakur2021beir} or other document retrieval benchmarks,
EABench data spans \emph{multiple modalities} (emails, meetings, files, chats,
channels) and requires \emph{cross-modal reasoning}.  Unlike curated enterprise
benchmarks such as WorkArena, EABench data can be generated on demand for any
domain and scale, enabling evaluation under distribution shift.

\subsection{Agent Configuration}
\label{sec:agentconfig}

EABench adopts a \emph{configuration-as-code} model in which every aspect of
an agent is specified in a human-readable YAML file.  This enables rigorous
ablation studies without code changes and allows practitioners to
version-control agent designs alongside benchmark data.

\textbf{AgentConfig schema.}
The top-level \texttt{AgentConfig} Pydantic model includes:

\begin{itemize}[leftmargin=*,itemsep=1pt]
  \item \texttt{model}: Provider type (\texttt{openai}, \texttt{azure},
    \texttt{anthropic}, \texttt{local}), model name, and sampling parameters
    (temperature, top-p, max tokens).
  \item \texttt{embedding}: Provider and model for dense retrieval.
  \item \texttt{system\_prompt}: The agent's instruction prompt, which may
    include a \texttt{\{user\_profile\}} placeholder automatically populated
    with the active user's profile at runtime.
  \item \texttt{planning\_prompt}: An optional prompt used exclusively during
    the planning phase of the Researcher strategy.
  \item \texttt{query\_analyzer\_prompt}: Per-tool LLM prompts for the query
    analyzer (see \S\ref{sec:tools}).
  \item \texttt{tools.definitions}: The list of tool names available to this
    agent, enabling fine-grained capability control.
  \item \texttt{flow.strategy}: Either \texttt{react} or \texttt{researcher}.
  \item \texttt{flow.max\_turns}: Maximum number of ReAct iterations before
    the agent is forced to return.
\end{itemize}

\textbf{Example ReAct agent configuration.}
Listing~\ref{lst:react_agent} shows a condensed YAML for a ReAct agent.
The \texttt{\{user\_profile\}} placeholder in \texttt{system\_prompt} is
resolved at runtime with the active user's profile, enabling per-user
personalisation without any code changes.
The \texttt{tools.definitions} list gives practitioners fine-grained
control over which tools are exposed to the model.

\begin{lstlisting}[caption={Condensed ReAct agent configuration (\texttt{react\_agent.yaml}).},label={lst:react_agent}]
id: react-agent-v1
version: "1.0.0"
model:
  provider: azure
  name: gpt-4o
  parameters: {temperature: 0.7}
embedding:
  provider: azure
  model: text-embedding-ada-002
system_prompt: |
  You are an intelligent enterprise assistant.
  ### User Context (Personalisation)
  {user_profile}   # resolved at runtime for the active user
  ### Tool Usage Guidelines
  - Search first: use search_email, search_file, search_meeting.
  - Resolve names to User IDs via search_people before other
    searches.
  - Use execute_python ONLY for data analysis; print all results.
  ### Citations
  - Cite every retrieved fact inline with [^N^] and add a
    References section at the end of every response.
tools:
  definitions:
    - search_email
    - search_file
    - search_meeting
    - search_chat
    - search_people
    - read_file
    - execute_python
flow:
  strategy: react
  max_turns: 5
\end{lstlisting}

\textbf{Example Plan-and-Execute (Researcher) agent configuration.}
Listing~\ref{lst:researcher_agent} shows the additional
\texttt{planning\_prompt} field used by the Researcher strategy.
The planner LLM receives the raw user query and the user profile, and
emits a structured step-by-step research plan that is then injected as
context for the standard ReAct execution loop.
Because planning and execution use separate model invocations, the
planner model can be a smaller, faster model (e.g., \texttt{gpt-4o-mini})
while the executor uses a more capable model.

\begin{lstlisting}[caption={Condensed Plan-and-Execute agent configuration (\texttt{researcher\_agent.yaml}).},label={lst:researcher_agent}]
id: researcher-agent-v1
version: "1.0.0"
model:
  provider: azure
  name: gpt-4o-mini
  parameters: {temperature: 0.0}
planning_prompt: |
  You are a Senior Research Planner.
  Create a step-by-step plan to answer the user's request.
  User Request: {user_query}
  User Profile:  {user_profile}
  Guidelines:
  1. Resolve all names to User IDs via search_people first.
  2. Break multi-hop questions into dependent steps.
  3. Do NOT execute -- output only the numbered plan.
  Output format:
  1. [Step Name]: description  (Tool: <tool_name>)
  2. [Step Name]: description  (Tool: <tool_name>)
system_prompt: |
  ... (same enterprise-assistant prompt as ReAct agent) ...
tools:
  definitions:
    - search_email
    - search_file
    - search_meeting
    - search_people
    - read_file
    - execute_python
flow:
  strategy: researcher
  max_turns: 10
\end{lstlisting}

\textbf{ReAct execution.}
Under the \texttt{react} strategy, the \texttt{AgentRunner} maintains a running
conversation history and iterates through LLM call $\to$ tool execution $\to$
observation cycles until the LLM produces a response without tool calls, or
\texttt{max\_turns} is reached.  The system prompt is injected once at the
beginning of the conversation; subsequent turns contain tool results as
\texttt{tool} role messages.

\textbf{Researcher (Plan-and-Execute) execution.}
Under the \texttt{researcher} strategy, \texttt{AgentRunner} first invokes the
LLM with the \texttt{planning\_prompt} to produce a high-level, step-by-step
research plan.  The plan is then injected into the conversation as context for
the standard ReAct loop.  This separation allows the planner to reason globally
about the task structure before committing to individual tool invocations.  The
final response is sanitized to remove internal planning artifacts before being
returned to the user.

\subsection{Enterprise Tool Suite}
\label{sec:tools}

EABench provides a comprehensive suite of enterprise tools organized into two
categories.

\textbf{Semantic search tools.}
Seven search tools cover every content type in the tenant:

\begin{itemize}[leftmargin=*,itemsep=1pt]
  \item \texttt{search\_email}: Multi-strategy email search supporting semantic
    similarity, keyword matching, sender filtering, and recency ranking.  A
    \emph{query analyzer} sub-module uses an LLM to determine the optimal
    retrieval strategy for each incoming query.
  \item \texttt{search\_file}: Vector-based search over indexed file content.
  \item \texttt{search\_meeting}: Search over meeting agendas and transcripts.
  \item \texttt{search\_chat} / \texttt{search\_group\_chat} /
    \texttt{search\_channel}: Search over 1-on-1 chats, group chats, and
    channel posts respectively.
  \item \texttt{search\_people}: User directory search by name, title, skill,
    or department.
  \item \texttt{search\_in\_file}: Keyword-based search within a specific file.
\end{itemize}

All search tools respect user-level access control: results are filtered to
documents accessible to the currently active user identity.  This enables
evaluation of realistic \emph{data visibility} scenarios where agents must
answer queries under specific permission contexts.

\textbf{Execution tools.}
Three tools enable agentic computation beyond retrieval:

\begin{itemize}[leftmargin=*,itemsep=1pt]
  \item \texttt{read\_file} / \texttt{list\_files}: Sandboxed file system access.
  \item \texttt{execute\_command}: Sandboxed shell command execution with
    timeout enforcement.
  \item \texttt{execute\_python}: Sandboxed Python code execution with
    stdout/stderr capture and new-file detection.
\end{itemize}

\textbf{Tool registration.}
New tools are registered using a \texttt{@registry.register} decorator, which
automatically generates a JSON Schema description from the Pydantic
\texttt{args\_schema} class.  Dependencies (\texttt{sandbox},
\texttt{search\_engine}, \texttt{llm}, \texttt{query\_analyzer}) are injected
automatically by the agent runner based on the tool function's parameter
signature, eliminating the need for manual dependency wiring.

\subsection{End-to-End Evaluation and User-Level Experience Simulation}
\label{sec:eval}

EABench's evaluation harness is designed to measure \emph{agent quality} across
multiple orthogonal dimensions rather than reducing performance to a single
score.

\textbf{Deterministic assertions.}
Each evaluation case includes a list of assertion descriptions---natural-language
statements about what the response should or should not contain.  At evaluation
time, a judge LLM assesses whether each assertion is satisfied, yielding a
discrete pass/fail result per assertion.  The \emph{assertion score} is the
fraction of passing assertions.

\textbf{Citation scoring.}
EABench adopts a structured citation format
(\texttt{[\^{}N\^{}] Title (Author, Date) [Type: <type>, ID: <id>]}) that
enables automated verification.  The evaluator extracts all citations from the
agent's response, fetches the corresponding source entities from the tenant,
and prompts a judge LLM to score the relevance of each citation to the user's
query.  This directly measures \emph{factual grounding}: an agent that
fabricates non-existent email IDs or cites an irrelevant meeting receives a low
citation score regardless of the plausibility of its narrative.

\textbf{LLM-as-a-Judge (side-by-side comparison).}
In addition to absolute scoring, EABench supports \emph{comparative evaluation}
of two agent configurations on the same set of queries.  A judge LLM receives
both responses and determines a winner based on the evaluation criteria,
producing a win-rate and a statistical significance estimate via a paired
$t$-test.  This enables ablation studies (e.g., does the Researcher strategy
outperform ReAct on multi-hop queries?) with quantified confidence.

\textbf{User-level experience simulation.}
A key feature of EABench is the ability to specify a \texttt{user\_id} per
evaluation case.  When a case is run under a specific user identity, the search
engine filters all retrieval results to documents accessible to that user
according to the tenant's group membership rules.  This simulates the
experience of different personas---an engineer, a manager, a contractor---and
enables evaluation of whether agents correctly respect data access boundaries.

\textbf{Metric aggregation.}
The final result for each evaluation case is determined by applying minimum
pass thresholds to both dimensions: a case \emph{passes} only when
$\text{assertion\_score} \geq 0.75$ \textbf{AND}
$\text{citation\_score} \geq 0.70$.
Neither threshold alone is sufficient, ensuring that an agent must both satisfy
the stated requirements and ground its response in real source entities.
Aggregate metrics across a full evaluation set include pass rate, mean
assertion score, mean citation score, mean latency, mean token usage, and mean
tool call count.

\textbf{Example evaluator configuration.}
Listing~\ref{lst:judge_config} shows the evaluator (judge) configuration.
The \texttt{prompts} block defines two LLM prompt templates:
\texttt{citation\_relevance} measures factual grounding by scoring whether
each cited source is genuinely relevant to the query, and
\texttt{assertion\_check} verifies whether the agent's response satisfies
each natural-language assertion, producing a per-assertion pass/fail result
and a weighted final score.

\begin{lstlisting}[caption={Condensed evaluator (judge) configuration (\texttt{default\_judge.yaml}).},label={lst:judge_config}]
name: "Default Judge Prompts"
model:
  provider: azure
  name: gpt-4o
  parameters: {temperature: 0.0}
prompts:
  citation_relevance: |
    You are an expert evaluator. Judge the relevance of citations
    used by the agent to answer the user query.
    User Query:       {query}
    Agent Tool Calls: {tool_calls}
    Agent Response:   {response}
    Score 0-1 (1 = perfectly relevant; 0 = hallucinated/irrelevant).
    Output YAML:
      score: <float 0.0-1.0>
      explanation: "<string>"
  assertion_check: |
    Verify whether the agent response satisfies each assertion.
    User Query:     {query}
    Agent Response: {response}
    Assertions:     {assertions}
    For each assertion output passed: true/false and reasoning.
    Compute final_score as the weighted fraction of passed assertions.
    Output YAML:
      assertions:
        - {id: 1, passed: <bool>, reasoning: "<string>"}
      final_score: <float 0.0-1.0>
      summary: "<string>"
\end{lstlisting}

\textbf{Example evaluation case.}
Listing~\ref{lst:eval_case} shows a representative evaluation case
produced by the automated case-generation pipeline for a tenant.
Each case binds a natural-language query to a specific user identity
(enforcing access control), a set of verifiable assertions, and a list
of expected source entities for citation checking.

\begin{lstlisting}[caption={Auto-generated evaluation case (TechNova tenant, \texttt{eval\_dataset.yaml}).},label={lst:eval_case}]
- id: case_001
  # Query run under Omar Khan's access context
  query: >
    Show me Priya Patel's email about client app
    performance issues.
  user_id: okhan
  assertions:
    - description: "Returns email titled 'Client Concerns About
                    App Performance'"
    - description: "Identifies sender as ppatel"
  entity_list:
    - entity_type: email
      entity_id: email-009
\end{lstlisting}

Table~\ref{tab:comparison} summarizes how EABench compares with existing
benchmarks across key dimensions.

\begin{table*}[t]
\centering
\caption{Comparison of EABench with existing agent benchmarks.
$\checkmark$ = supported, -- = not supported.}
\label{tab:comparison}
\small
\begin{tabular}{lcccccc}
\toprule
\textbf{Feature} & \textbf{WebArena} & \textbf{AgentBench} & \textbf{ToolBench} & \textbf{$\tau$-bench} & \textbf{WorkArena} & \textbf{EABench} \\
\midrule
Multi-source enterprise data          & --            & --             & --             & --              & Partial        & $\checkmark$ \\
Synthetic data generation             & --            & --             & --             & --              & --             & $\checkmark$ \\
Configurable agent strategies         & --            & --             & --             & --              & --             & $\checkmark$ \\
User-level access control             & --            & --             & --             & --              & --             & $\checkmark$ \\
LLM-as-a-Judge evaluation             & --            & --             & --             & --              & --             & $\checkmark$ \\
Citation grounding verification       & --            & --             & --             & --              & --             & $\checkmark$ \\
Side-by-side comparison               & --            & --             & --             & --              & --             & $\checkmark$ \\
Multi-provider LLM support            & --            & Partial        & --             & --              & --             & $\checkmark$ \\
\bottomrule
\end{tabular}
\end{table*}

% ══════════════════════════════════════════════════════════════════════
% 4  EXPERIMENTS
% ══════════════════════════════════════════════════════════════════════
\section{Experiments}
\label{sec:experiments}

This section is reserved for experimental results.  Forthcoming work will
report benchmark scores across multiple LLM backends (GPT-4o, Claude 3.5
Sonnet, Gemini 1.5 Pro, open-source models) and agent configurations (ReAct
vs.\ Researcher, varying tool subsets, varying prompt designs) on a
standardized collection of enterprise tenants spanning at least three industry
verticals (technology, healthcare, finance).  Experiments will investigate:

\begin{itemize}[leftmargin=*,itemsep=2pt]
  \item How agent strategy (ReAct vs.\ Researcher) affects performance as a
    function of query complexity (single-hop vs.\ multi-hop).
  \item The impact of LLM backbone on citation accuracy and assertion pass
    rate.
  \item The effect of available tool subsets on task success rate.
  \item Sensitivity of evaluation metrics to judge LLM quality.
  \item Scaling behavior of the data generator (small: 5 users, 7 days vs.\
    large: 50 users, 90 days).
\end{itemize}

% ══════════════════════════════════════════════════════════════════════
% 5  CONCLUSION AND DISCUSSION
% ══════════════════════════════════════════════════════════════════════
\section{Conclusion and Discussion}
\label{sec:conclusion}

We have presented \textbf{EABench}, a configurable benchmark framework for
evaluating LLM-powered agents in enterprise settings.  EABench addresses three
gaps in existing evaluation infrastructure: the absence of realistic, at-scale
enterprise corpora; the lack of flexible, code-free agent architecture
configuration; and the insufficiency of single-metric evaluation for multi-step
agentic systems.

\textbf{Key contributions.}
EABench's LLM-driven data generator produces coherent, inter-connected
organizational data across six content types (emails, meetings, files, chats,
group chats, channels) with controllable scale and narrative structure.  Its
configuration-as-code agent runtime enables systematic comparison of ReAct and
Plan-and-Execute strategies across arbitrary LLM backends without modifying
source code.  Its multi-dimensional evaluation harness---combining deterministic
assertions, citation grounding verification, and LLM-as-a-Judge---provides a
more complete picture of agent quality than any single metric alone.
User-level experience simulation adds a privacy and access-control dimension
that is absent from prior benchmarks.

\textbf{Limitations.}
EABench currently supports only English-language content generation and
English-language agent evaluation.  The data generator relies on a capable LLM
(GPT-4-class or better) to produce coherent narratives; lower-quality models
may produce less realistic data.  The citation format required by the evaluator
must be enforced through the system prompt; agents that do not adhere to this
format receive a zero citation score even if their responses are factually
correct.  Finally, the evaluation assertions are LLM-generated and may not
capture all dimensions of response quality.

\textbf{Future work.}
Several directions are planned:

\begin{enumerate}[leftmargin=*,itemsep=2pt]
  \item \emph{Docker sandbox integration}: Extend the sandbox backend to
    support containerized execution for stronger security guarantees when
    evaluating agents that run arbitrary code.
  \item \emph{Adversarial tenant generation}: Create tenants with deliberate
    ambiguities, contradictory documents, and planted misinformation to
    evaluate agent robustness.
  \item \emph{Multi-agent evaluation}: Extend the runtime to support
    multi-agent workflows, where specialized sub-agents (e.g., an
    email-reading agent and a meeting-summarizing agent) collaborate on a
    shared task.
  \item \emph{Multilingual support}: Extend the data generator and evaluation
    harness to support non-English enterprise environments.
  \item \emph{Human evaluation calibration}: Conduct a study comparing
    LLM-as-a-Judge scores with human annotations to quantify judge
    reliability.
\end{enumerate}

We believe EABench fills a critical gap between the rapid development of LLM
agent capabilities and the infrastructure needed to evaluate them rigorously in
the settings where they will actually be deployed.  We release all code, data
generation scripts, and evaluation tooling under the MIT license to support
reproducible research in enterprise AI.

% ══════════════════════════════════════════════════════════════════════
% REFERENCES
% ══════════════════════════════════════════════════════════════════════
\bibliographystyle{plainnat}

\begin{thebibliography}{15}

\bibitem[Chiang et~al.(2024)]{chiang2024chatbot}
W.-L. Chiang, L.~Zheng, Y.~Sheng, A.~N.~Angelopoulos, T.~Li, D.~Li, et~al.
\newblock Chatbot Arena: An open platform for evaluating LLMs by human
  preference.
\newblock \emph{arXiv preprint arXiv:2403.04132}, 2024.

\bibitem[Drouin et~al.(2024)]{drouin2024workarena}
A.~Drouin, M.~Gasse, M.~Caccia, I.~H.~Laradji, M.~Del~Verme, T.~Marty,
  et~al.
\newblock WorkArena: How capable are web agents at solving common knowledge
  work tasks?
\newblock \emph{arXiv preprint arXiv:2403.07718}, 2024.

\bibitem[Glazer et~al.(2024)]{glazer2024frontiermath}
E.~Glazer, E.~Erdil, T.~Besiroglu, D.~Chicharro, E.~Chen, A.~Gunning,
  et~al.
\newblock FrontierMath: A benchmark for evaluating advanced mathematical
  reasoning in AI.
\newblock \emph{arXiv preprint arXiv:2411.04872}, 2024.

\bibitem[Liu et~al.(2023)]{liu2023agentbench}
X.~Liu, H.~Yu, H.~Zhang, Y.~Xu, X.~Lei, H.~Lai, et~al.
\newblock AgentBench: Evaluating LLMs as agents.
\newblock \emph{arXiv preprint arXiv:2308.03688}, 2023.

\bibitem[Microsoft(2023)]{microsoft2023copilot}
Microsoft.
\newblock Microsoft 365 Copilot: The AI-powered productivity tool for work.
\newblock Microsoft Blog, 2023.

\bibitem[Mitra et~al.(2024)]{mitra2024agentinstruct}
A.~Mitra, L.~Del~Corro, S.~Mahajan, A.~Codas, C.~Simoes, S.~Agrawal,
  et~al.
\newblock AgentInstruct: Toward generative teaching with agentic flows.
\newblock \emph{arXiv preprint arXiv:2407.03502}, 2024.

\bibitem[Qin et~al.(2023)]{qin2023toolllm}
Y.~Qin, S.~Liang, Y.~Ye, K.~Zhu, L.~Yan, Y.~Lu, et~al.
\newblock ToolLLM: Facilitating large language models to master 16000+
  real-world APIs.
\newblock \emph{arXiv preprint arXiv:2307.16789}, 2023.

\bibitem[Thakur et~al.(2021)]{thakur2021beir}
N.~Thakur, N.~Reimers, A.~R{\"u}ckl{\'e}, A.~Srivastava, and I.~Gurevych.
\newblock BEIR: A heterogeneous benchmark for zero-shot evaluation of
  information retrieval models.
\newblock In \emph{NeurIPS 2021 Datasets and Benchmarks}, 2021.

\bibitem[Wang et~al.(2023)]{wang2023plansolve}
L.~Wang, W.~Xu, Y.~Lan, Z.~Hu, Y.~Lan, R.~K.-W.~Lee, and E.-P.~Lim.
\newblock Plan-and-Solve prompting: Improving zero-shot chain-of-thought
  reasoning by large language models.
\newblock In \emph{ACL}, 2023.

\bibitem[Yao et~al.(2022)]{yao2022react}
S.~Yao, J.~Zhao, D.~Yu, N.~Du, I.~Shafran, K.~Narasimhan, and Y.~Cao.
\newblock ReAct: Synergizing reasoning and acting in language models.
\newblock \emph{arXiv preprint arXiv:2210.03629}, 2022.

\bibitem[Yao et~al.(2024)]{yao2024taubench}
S.~Yao, D.~Yu, J.~Zhao, I.~Shafran, T.~L.~Griffiths, Y.~Cao, and
  K.~Narasimhan.
\newblock $\tau$-bench: A benchmark for tool-agent-user interaction in
  real-world domains.
\newblock \emph{arXiv preprint arXiv:2406.12045}, 2024.

\bibitem[Zhao et~al.(2024)]{zhao2024wildchat}
W.~Zhao, X.~Ren, J.~Hessel, C.~Cardie, Y.~Choi, and Y.~Deng.
\newblock WildChat: 1M ChatGPT interaction logs in the wild.
\newblock In \emph{ICLR}, 2024.

\bibitem[Zheng et~al.(2023)]{zheng2023judging}
L.~Zheng, W.-L.~Chiang, Y.~Sheng, S.~Zhuang, Z.~Wu, Y.~Zhuang, et~al.
\newblock Judging LLM-as-a-judge with MT-Bench and Chatbot Arena.
\newblock In \emph{NeurIPS}, 2023.

\bibitem[Zhou et~al.(2023)]{zhou2023webarena}
S.~Zhou, F.~F.~Xu, H.~Zhu, X.~Zhou, R.~Lo, S.~Sridhar, et~al.
\newblock WebArena: A realistic web environment for building autonomous agents.
\newblock \emph{arXiv preprint arXiv:2307.13854}, 2023.

\end{thebibliography}

\end{document}
